\subsection{Van Kampen's System Size Expansion}
\begin{enumerate}
\item The system size expansion is a technique, pioneered by van Kampen  \citep{vankampen2011}, to find an approximation to the solution of the master equation of a stochastic system. Within the so-called linear noise approximation, the method yields an equation for the temporal evolution of the deterministic state variable (Eq. (\ref{linear_noise_D})), and an equation for the probability density of the fluctuations around the deterministic prediction (Eq. (\ref{linear_noise_F})):
\begin{eqnarray}
\label{linear_noise_D}
\frac{d\Phi}{dt}                & = &   \alpha_1(\Phi)  \\
\frac{\partial\Pi}{\partial\,t} & = & - \frac{d\alpha_1}{d\Phi}\frac{\partial(\xi\Pi)}{\partial\,\xi}+\frac{1}{2}\alpha_2(\Phi)\frac{\partial^2\Pi}{\partial\,\xi^2}
\label{linear_noise_F}
\end{eqnarray}
where coefficients $\alpha'_1(\Phi)$ and $\alpha_2(\Phi)$ can be obtained from the {\em macroscopic law}\footnote{See Eq. (\ref{linear_noise_D}) and notes of section 2 for details.}. Check that Eq. (\ref{linear_noise_F}) admits as a solution the following probability density of Gaussian type:
\begin{equation}
\label{Piofxit}
\Pi(\xi,t) = \frac{1}{\sqrt{2\pi\,<\xi^2>}}\,\exp \left(-\frac{\xi^2}{2\,<\xi^2>}\right)
\end{equation}
provided that $<\xi> = 0$, and that the second moment, $<\xi^2>$, evolves in time according to the following ODE:
\begin{equation}
\partial_t<\xi^2> = 2\,\frac{d\alpha_1}{d\Phi}\,<\xi^2> \,+\, \alpha_2 
\label{eq:xi}
\end{equation}

\item Consider individuals arriving to an island where individuals do not reproduce. In this case, we say that this island acts as a {\em sink} ---as opposite to a {\em source}--- for this population. If we define the parameter $\lambda_0$ as the island immigration rate per unit area, then the probability for the island population to increase by one in a short time interval, $dt$, will be proportional to island area:
\begin{equation}
  \Pb(n \rightarrow n + 1) = \lambda_0 A\, dt 
\end{equation} 
If we assume that per capita mortality rates are density-independent, then the probability for the island population to decrease by one in a short time interval, $dt$, will be:  
\begin{equation}
  \Pb(n \rightarrow n - 1) = \mu \,n\, dt 
\end{equation}  
where $n$ is the number of individuals in the island at a give time, $t$. In this system, address the follwing questions: 
\begin{itemize} 
\item Using the operators $E$ and $E^{-1}$, write the master equation. Solve the equation for the temporal evolution of the density of individuals in the island\footnote{See also problem 4(b)iii in homework 2.}, $\Phi$ ---the {\em macroscopic law}. 
\item Under the linear-noise approximation, solve for the temporal evolution of the variance of the fluctuations,  $<\xi^2>$, around the deterministic density ---see Eq. (\ref{eq:xi}). Assess to what extend the Gaussian probability:
\begin{equation}
  P(n,t) = \frac{1}{\sqrt{2\pi\,A\,\textup{Var}(\xi)}}\exp\left(-\frac{(n-A\Phi(t))^2}{2\,A\,\textup{Var}(\xi)}\right)
\end{equation} 
is a good approximation to the exact probability of finding $n$ individuals in the island at time $t$, $P(n,t)$. Take the empty island as the initial condition. The exact probability distribution at increasing times can be calculated through simulation using Gillespie algorithm. 
\end{itemize}

\item Consider a tiny reactor of volume $V$ containing a solution of chemical species $A^{\star}$ and $A$. $A^{\star}$ is a stable precursor of chemical species $A$. Chemical species $A^{\star}$ converts into $A$ through the following auto-catalytic, non-reversible reaction:
\begin{equation}
A^{\star} + A \longrightarrow A + A
\end{equation}    
However, chemical species $A$ has a tendency to form non-soluble dimers forming precipitates. The precipitation reaction can be represented by:   
\begin{equation}
A + A \longrightarrow \emptyset
\end{equation}
because random encounters of two $A$ molecules involve their disappearance from the solution. Consider that the concentration of the precursor $A^{\star}$ is externanlly controlled to a constant level, $[A^{\star}] = \Phi^{\star}$, and that the concentration of the chemical species $A$ at the initial time is $[A]= \Phi_A(0) = \Phi_0$. If $n$ is the number of molecules of species $A$ in the solution, describe the stochastic dynamics of this system by addressing the following questions:
\begin{itemize}
\item Write the probability rates $g_n$ and $r_n$ of the number of molecules of species $A$ increasing and decreasing by 1 per unit time,  respectively. Using the operators $E$ and $E^{-1}$, write the resulting master equation for the temporal evolution of the probability of having $n$ molecules of species $A$ in solution at any time $t$, $P(n,t)$. 
\item Write and solve the equation for the temporal evolution of the concentration of species $A$ ---the {\em macroscopic law}.    
\item Under the linear-noise approximation, solve for the temporal evolution of the variance of the fluctuations,  $<\xi^2>$, around the deterministic density ---see Eq. (\ref{eq:xi}). Assess to what extend the Gaussian probability:
\begin{equation}
  P(n,t) = \frac{1}{\sqrt{2\pi\,V\,\textup{Var}(\xi)}}\exp\left(-\frac{(n-V\Phi(t))^2}{2\,V\,\textup{Var}(\xi)}\right)
\end{equation} 
is a good approximation to the exact probability of finding $n$ molecules in the reactor at time $t$, $P(n,t)$. Consider, initially, there are $n_0$ ($n_0 = \Phi_0\,V$) molecules of chemical species, $A$, this is, $P(n_0, 0) = 1$. The exact probability distribution at increasing times can be calculated through numerical integration of the set of ODEs defining the master equation. 
\end{itemize}
\end{enumerate}

